\documentclass[14pt,a4paper]{extarticle}

\usepackage[utf8]{inputenc}
\usepackage[T1]{fontenc}
\usepackage[english]{babel}
\usepackage{fontspec}
\setmainfont{Times New Roman}

\usepackage[
    top=2cm,
    bottom=2cm,
    left=3cm,
    right=1.5cm
]{geometry}

\usepackage{setspace}
\usepackage{indentfirst}
\usepackage{enumitem}
\usepackage{booktabs}
\usepackage{longtable}
\usepackage{array}
\usepackage{hyperref}
\usepackage{xcolor}
\usepackage{titlesec}
\usepackage{caption}
\usepackage{tabularx}
\usepackage{amsmath}
\usepackage{float}
\usepackage{makecell}

\newcolumntype{L}[1]{>{\raggedright\arraybackslash}p{#1}}

\hypersetup{
    colorlinks=true,
    linkcolor=black,
    urlcolor=blue,
    citecolor=black
}

\onehalfspacing
\setlength{\parindent}{1.25cm}
\setlength{\parskip}{0pt}

\setlist[itemize]{leftmargin=1.5em}
\setlist[enumerate]{leftmargin=1.8em}

\titleformat{\section}{\large\bfseries}{\thesection.}{0.5em}{}
\titleformat{\subsection}{\normalsize\bfseries}{\thesubsection.}{0.5em}{}
\titleformat{\subsubsection}{\normalsize\itshape}{\thesubsubsection.}{0.5em}{}

\begin{document}

\begin{titlepage}
\begin{center}

{\normalsize
Federal Autonomous Educational Institution of Higher Professional Education\\
``National Research University Higher School of Economics''
}

\vspace{1.2cm}

{\normalsize
Faculty of Social Sciences\\
Educational Programme ``Computational Social Sciences''
}

\vspace{2.4cm}

{\normalsize
Course work on the topic:
}

\vspace{0.6cm}

{\large\bfseries
``Synthetic sample generation in the context of experimental research''
}

\vspace{1cm}

\vfill

\begin{flushright}
\begin{minipage}{0.48\textwidth}
\raggedright
\textbf{Scientific advisor:}\\
Evgeny Sedashov,\\
Programme Academic Supervisor ``Computational Social Sciences'',\\
PhD, Binghamton, \\
State University of New York 
\end{minipage}
\end{flushright}

\vspace{1cm}

\begin{flushright}
\begin{minipage}{0.48\textwidth}
\raggedright
\textbf{Completed by:}\\
Maria Akhryapina,\\
Elizaveta Kuznechenkova,\\
Maria Nesterskaya\\
3rd-year students of group BVSN231
\end{minipage}
\end{flushright}

\vfill

{\normalsize
Moscow, 2026
}

\end{center}
\end{titlepage}

\tableofcontents
\newpage

\section{Introduction}

Large language models (LLMs) are increasingly discussed as tools for generating synthetic respondents, stress-testing questionnaires, and piloting experimental designs before launching costly human data collection. In this project, LLMs are treated not as replacements for survey respondents in a substantive sense, but as instruments for \emph{synthetic sample generation} in the context of early-stage experimental research. The practical motivation is straightforward: if a model can approximate how different kinds of respondents react to alternative wordings, vignettes, and contextual cues, it may be useful for piloting the structure of a survey experiment, identifying fragile manipulations, and revealing where a design is likely to be distorted by response biases.

The central research question of this study is the following: \textbf{How do large language models react to stimuli in vignette-based survey experiments?} This question is narrower than the broader agenda of ``LLMs as synthetic respondents'' and is intentionally focused on one concrete challenge: whether prompt-level interventions can shift the model's response pattern in morally sensitive survey settings, and whether this shift is visible both in the final answer and in the reasoning trace that accompanies it.

The project is situated in the domain of \emph{social desirability bias} (SDB). This choice is theoretically and methodologically consequential. SDB-related survey items are especially demanding for LLM-based piloting because they lie at the intersection of two response-shaping forces. First, human respondents often adjust their answers toward what is socially acceptable, normatively safe, or reputationally advantageous. Second, contemporary LLMs are themselves tuned to avoid harmful, offensive, or normatively problematic outputs. As a result, SDB is not merely one topic among many. It is a difficult test case in which the model's safety alignment may interact with the substantive logic of the survey instrument itself.

This makes SDB particularly valuable for studying whether LLMs can serve as usable pilot environments for vignette experiments. If the model reproduces only sanitized, ``model citizen'' answers regardless of treatment, it is of little use as a tool for experimental pretesting in sensitive domains. If, however, it reacts systematically to experimental manipulations and carefully designed prompt stimuli, then it may help researchers understand whether a survey design can elicit meaningful variation before the design is fielded on a human sample.

The present text is a design-oriented coursework draft rather than a final empirical article. It does not present completed estimates or final conclusions. Instead, it introduces the logic of the project, motivates the selection of the SDB domain, describes the initial experimental corpus, explains the model and prompting setup, outlines the current system-prompt stimuli, and presents two planned analytic tracks: structured analysis of reasoning traces and layer-wise analysis of hidden states.

\section{Why Social Desirability Bias is a Demanding Domain for LLM-Based Piloting}

For the purposes of this project, survey experiments related to social desirability bias can be divided into three broad categories.

\begin{enumerate}
    \item \textbf{Methodological SDB experiments.} These studies directly test techniques for reducing, detecting, or re-framing socially desirable responding. They are often the most difficult to translate into an LLM setting because the treatment is aimed at the response process itself rather than solely at the substantive content of the question.
    \item \textbf{Self-report SDB experiments.} These studies ask respondents to report their own behavior, intentions, or personal characteristics in sensitive domains. They are important but harder to simulate cleanly, because the prompt must induce both a plausible self-concept and a plausible response style.
    \item \textbf{Attitude-expression SDB experiments.} These studies ask respondents to express an evaluative judgment, moral opinion, or normative stance toward a vignette, policy, or action. They preserve the core challenge of socially desirable responding while avoiding some of the additional complexity of autobiographical self-report.
\end{enumerate}

This project starts with the third category, namely attitude-expression experiments. There are both conceptual and practical reasons for doing so. Conceptually, attitude-expression designs make it easier to isolate how the model reacts to a vignette manipulation itself: the response is usually a judgment about an externally described case rather than a self-description. Practically, they are well suited to a first-stage pilot because many published attitude experiments have compact designs, clearly differentiated treatment conditions, and outcome measures that can be mapped onto a stable prompt template.

At the same time, attitude-expression experiments remain genuinely difficult for LLMs. When the vignette concerns violence, punishment, prejudice, civic duty, or generosity, the model may not simply simulate a respondent. It may instead drift toward a normatively polished answer that looks maximally safe. This is exactly why the SDB domain is useful here: it offers a setting in which the model's alignment and the survey's substantive logic are likely to collide.

\section{Substantive Focus: Attitude-Expression Experiments}

\subsection{Why the project begins with attitude expression}

The long-term research agenda is broader than the present draft. In principle, the project is interested in all three categories listed above. However, the current empirical stage focuses only on attitude-expression experiments. This restriction is intentional. It narrows the design to a tractable subset while preserving the central methodological difficulty of interest, namely whether an LLM can react to morally loaded vignettes in a way that is neither trivially sanitized nor artificially provocative.

Attitude-expression experiments are especially useful for an initial study because they allow the researcher to compare three levels at once: the original treatment embedded in the vignette, the model's final answer to the survey question, and the reasoning trace through which the answer is reached. This makes them a suitable entry point for studying both observable response patterns and internal argumentative dynamics.

\subsection{Case selection logic}

The current corpus is deliberately small and functions as a pilot subset. The two initial studies were selected according to the following criteria:

\begin{itemize}
    \item the article contains, or allows a defensible reconstruction of, the experimental conditions;
    \item the design includes multiple conditions so that treatment contrasts can be meaningfully reproduced;
    \item the main dependent variable is an evaluative judgment or normative stance rather than autobiographical self-report;
    \item the wording can be implemented in a relatively stable prompt format across models;
    \item the selected set spans both a socially \emph{negative} and a socially \emph{positive} direction of desirability pressure.
\end{itemize}

The last criterion is central. The present corpus was intentionally built around one case where the socially desirable tendency is likely to push toward condemnation of a harmful act, and one case where it may push toward overstatement of a prosocial act. This contrast makes it possible to study whether prompt stimuli operate symmetrically across different normative valences.

\subsection{Study 1: coercive interrogation / torture vignette}

The first study is based on a torture-related vignette derived from Hassner's work on critiques of torture. In the current implementation, the respondent is asked to evaluate whether a terror suspect should be subjected to sleep deprivation during interrogation. The experimental conditions vary the expected informational value of the coercive act and the severity or duration of the procedure.

This case serves as the \emph{negative} pole of the current corpus. Supporting coercive interrogation is normatively problematic, and a default socially desirable response is likely to move toward categorical rejection. At the same time, the treatment details make the vignette nontrivial: the model can appeal to consequences, necessity, proportionality, or moral prohibition. This makes the case useful for examining whether prompt stimuli merely change the extremity of condemnation or whether they alter how the model balances competing moral frames.

\subsection{Study 2: charitable donation vignette}

The second study is based on a donation experiment in which respondents evaluate willingness to donate under varying social conditions. In the current implementation, the conditions contrast a baseline private donation setting with settings in which the donor's act may be publicly recognized or socially visible.

This case serves as the \emph{positive} pole of the corpus. Donating to charity is a socially approved action, and the socially desirable tendency here may work in the opposite direction: rather than hiding an undesirable preference, the model may overstate its willingness to perform a virtuous act. For that reason, the donation case is analytically complementary to the torture case. Together, the two experiments allow the project to ask whether the same prompt stimulus reduces sanitization in morally opposite settings, or whether distinct mechanisms are at work in ``negative'' and ``positive'' SDB domains.

\subsection{Current corpus summary}

\begin{table}[H]
\centering
\caption{Current pilot corpus of attitude-expression experiments}
\begin{tabular}{p{3.6cm}p{3.2cm}p{7cm}}
\toprule
\textbf{Study} & \textbf{Normative direction} & \textbf{Role in current design} \\
\midrule
Torture / coercive interrogation vignette & Negative case & Tests whether prompt stimuli reduce safe condemnation and change how the model reasons about a morally problematic act \\
Charitable donation vignette & Positive case & Tests whether prompt stimuli reduce virtue-signaling or overstatement of prosocial willingness \\
\bottomrule
\end{tabular}
\end{table}

\section{Model Set and Generation Protocol}

\subsection{Why reasoning-capable models are used}

The current pilot is run on four models that expose or support reasoning-style generation: \texttt{moonshotai/kimi-k2-thinking}, \texttt{qwen3-9b}, \texttt{glm-4.6}, and \texttt{minimax-m2.5}. These models were selected for a pragmatic reason rather than because they represent an exhaustive benchmark. The project needs systems that are capable of generating not only a final answer but also a sufficiently rich reasoning trace that can later be analyzed as text and, in a separate analytic track, inspected through hidden-state representations.

The selected models also provide useful diversity. They differ in family, scale, alignment profile, and style of reasoning output. This heterogeneity is analytically valuable because one of the project's goals is to compare whether similar survey answers are produced through different argumentative patterns. A narrow one-family benchmark would be much less informative for that purpose.

\subsection{Deterministic setting and reproducibility}

All main generations are run at temperature $0.00$. The purpose of this choice is not to claim that deterministic decoding is behaviorally ``natural,'' but to maximize reproducibility during the piloting stage. Since the project compares subtle differences induced by vignette treatments and system-prompt stimuli, it is important to reduce avoidable generation noise and make the effect of the manipulated inputs easier to interpret.

This is especially important because the design is already high-dimensional: answers may vary by experiment, treatment condition, system-prompt stimulus, and model family. Fixing temperature at zero makes the comparison sharper and allows the researcher to interpret observed differences as more plausibly related to the manipulations rather than to uncontrolled sampling variance.

\subsection{Single-pass collection of outputs}

The study does \emph{not} use a two-stage or sequential setup. Instead, each prompt is run in a single pass, and from that single generation the project collects two outputs simultaneously:

\begin{enumerate}
    \item the \textbf{final answer} required by the survey instruction;
    \item the \textbf{reasoning trace} produced by the model while arriving at that answer.
\end{enumerate}

This distinction matters. The project is not comparing two separate tasks or two independent calls to the model. Rather, it treats the final answer and the reasoning trace as two analytically distinct components of the same response event. This preserves the tightest possible link between the stated outcome and the internal argumentative path available for inspection.

\section{System-Prompt Stimuli}

\subsection{Why stimulus design is central to the project}

The main experimental intervention introduced by the project is not only the original vignette treatment inherited from the source studies, but also a set of additional \emph{system-prompt stimuli}. These stimuli are designed to change how the model approaches the survey task itself. Their shared purpose is to reduce overly sanitized answering and to push the model toward what the project calls \emph{behavioral realism}: an attempt to approximate how a plausible respondent would actually answer under the described conditions, rather than how an idealized moral agent ought to answer.

The logic behind this family of stimuli is grounded in the methodological problem identified earlier. In SDB-related settings, a model may respond in a way that is safe, cautious, and normatively polished, but methodologically unhelpful. Such a pattern is undesirable if the goal is to pilot a survey instrument that is meant to capture human variation, including variation that is normatively uncomfortable. The system-prompt interventions are therefore designed as attempts to reframe the model's task from self-presentation toward measurement.

\subsection{Stimulus dimensions}

The current set of stimuli varies along three dimensions.

\begin{enumerate}
    \item \textbf{Harshness or force of the instruction.} Some stimuli directly warn about losses or penalties if responses look sanitized, while others formulate the realism request more cautiously or balance it with an anti-overcorrection clause.
    \item \textbf{Recipient of the consequences.} The stimuli vary who is described as suffering if the model produces unrealistic responses: a researcher, science as a public enterprise, a company, the public, or users of a product.
    \item \textbf{Valence of the framing.} Some stimuli use a negative frame (bad consequences if the answers are unrealistic), while others use a positive frame (benefits if the answers are realistic).
\end{enumerate}

These dimensions were not chosen arbitrarily. They map onto plausible mechanisms through which a system prompt may influence an aligned model. A harsher prompt may increase compliance but also provoke overcorrection. A prompt framed around science or public good may activate a different justification structure than one framed around company losses. A positive frame may be easier for the model to adopt without becoming defensive or excessively literal. In other words, the stimuli are not only interventions; they are miniature theories about what kinds of meta-instruction matter for response generation.

\subsection{Current stimulus inventory}

Table~\ref{tab:stimuli} summarizes the current inventory.

\begin{longtable}{L{4.5cm} L{8.5cm}}
\caption{Current system-prompt stimuli}\label{tab:stimuli}\\
\toprule
\textbf{Name} & \textbf{Core logic} \\
\midrule
\endfirsthead

\toprule
\textbf{Name} & \textbf{Core logic} \\
\midrule
\endhead

\makecell[l]{A\_original\_core\_\\realism} &
Strong realism instruction with negative consequences for the researcher if the responses look sanitized rather than human. \\

\makecell[l]{C\_balanced\_\\realism} &
Same core realism instruction, but with an explicit warning not to overcorrect toward provocative or implausibly extreme answers. \\

\makecell[l]{H\_science\_\\personal\_gain} &
Negative framing centered on scientific work suffering if unrealistic answers reduce the informativeness of the study; includes an anti-overcorrection clause. \\

\makecell[l]{I\_science\_\\public\_good} &
Negative framing centered on lower scientific value for broader public understanding; includes an anti-overcorrection clause. \\

\makecell[l]{J\_company\_\\personal\_gain} &
Negative framing centered on firm-level losses if misleading data distort product design; includes an anti-overcorrection clause. \\

\makecell[l]{K\_company\_\\public\_good} &
Negative framing centered on reduced usefulness of the product for real people; includes an anti-overcorrection clause. \\

\makecell[l]{L\_science\_public\_\\good\_positive} &
Positive framing: realistic responses improve scientific value and produce a more truthful picture of human behavior. \\

\makecell[l]{M\_company\_personal\_\\gain\_positive} &
Positive framing: realistic responses improve product design and benefit the company. \\

\makecell[l]{X\_measurement\_\\neutrality} &
Explicit measurement frame: the task is not to be morally admirable, cautious, or provocative, but to approximate a plausible respondent under the described conditions. \\

\bottomrule
\end{longtable}

\subsection{Interpretation of the stimulus families}

The nine current stimuli can be grouped into several analytically meaningful families.

\paragraph{1. Core realism and balanced realism.}
The baseline realism family consists of the original strong instruction (A) and its balanced variant (C). Their shared purpose is to directly challenge a sanitized ``model citizen'' response pattern. The balanced version adds a crucial methodological safeguard: it explicitly instructs the model not to become artificially contrarian or extreme. This family is meant to test the trade-off between stimulus force and response plausibility.

\paragraph{2. Recipient-of-consequence variants.}
The science/company variants (H, I, J, K) hold the core realism logic constant while varying who suffers if the model produces unrealistic answers and how that harm is framed. Some versions emphasize damage to the researcher's work or the company's value; others emphasize reduced public understanding or lower usefulness to real people. This family is designed to test whether the model is more responsive to personal, institutional, or public-good consequences.

\paragraph{3. Positive-framing variants.}
The positive versions (L, M) preserve the realism goal but switch from a penalty frame to a benefit frame. Instead of warning that unrealistic answers will produce harm, they state that realistic answers will improve science or product quality. This family is motivated by the possibility that a positively framed instruction may reduce defensiveness and elicit more natural departures from socially desirable responding.

\paragraph{4. Measurement neutrality.}
The neutrality stimulus (X) is conceptually distinct from the realism-loss family. Rather than emphasizing harm or benefit, it reframes the task itself: not moral self-presentation, not caution, not provocation, but measurement. This stimulus is particularly important because it comes closest to a methodological meta-instruction that one might plausibly want in a survey-piloting workflow.

\subsection{Current preliminary observations}

At the current stage, only preliminary tendencies can be stated.

\begin{itemize}
    \item \textbf{Harsh unbalanced stimuli appear to make answers sharper.} This is currently a tentative observation rather than a stable result. The word ``sharper'' still requires analytical refinement: it may refer to greater extremity of final answers, stronger normative language in the reasoning trace, lower expressed uncertainty, or more explicit rejection of safe-answer hedging.
    \item \textbf{Positive framing currently appears more effective at eliciting socially undesirable answers.} This finding is especially interesting because it suggests that moving away from punishment language may sometimes produce stronger departures from sanitized output than a direct warning does.
    \item \textbf{The identity of the harmed recipient does not currently seem to matter in the negative framing setup.} In other words, whether the consequences are described in terms of science, company loss, personal cost, or public value has not yet produced a clear difference in the negative versions. The same question remains open for the positive variants.
\end{itemize}

These observations are best treated as working hypotheses for the next stage of analysis rather than as completed empirical claims.

\section{Reasoning Trace Annotation and the MoralLens Instrument}

\subsection{Why reasoning traces require a dedicated annotation instrument}

Because the project collects reasoning traces together with final answers, it needs a principled way to transform long free-form text into a structured analytical representation. Raw reasoning traces are rich but unwieldy. They may contain explicit moral arguments, hedging, meta-commentary on the survey format, treatment uptake, stimulus uptake, or artifacts of alignment behavior. A dedicated annotation instrument is therefore required if these texts are to be analyzed systematically rather than impressionistically.

The annotation component of the project is being developed under the working name \textbf{MoralLens}. Its purpose is to take a model's reasoning trace and convert it into a structured set of variables describing how the model framed the task, what kinds of arguments it used, how it related to the treatment, and whether it displayed signs of social-desirability management or role leakage. According to the project presentation, MoralLens is conceived specifically as an automated analyzer of moral reasoning in LLM outputs, with the final goal of producing a structured summary of how the prompt stimuli affected the model's internal decision path.\footnote{Internal project presentation ``MoralLens: Automated analysis of moral reasoning in language models.''}

\subsection{Current analytical work before model training}

Before training the annotation model itself, the project requires a separate stage of analytic work devoted to the construction of a coding scheme. At the moment, manual inspection of the reasoning traces suggests that the relevant features fall into two broad groups.

\begin{enumerate}
    \item \textbf{Format-related features.} These are reflections on the survey form, answer scale, respondent role, or task interpretation. Examples include reasoning such as ``category 5 would look too saintly,'' ``the question asks what the respondent would think rather than what they would do,'' or explicit calibration of what counts as a plausible versus performatively moral answer.
    \item \textbf{Substantive features.} These concern the content of the vignette and the moral categories used to process it: where the treatment is located in the scenario, what normative or consequentialist principles the model invokes, how it justifies approval or disapproval, and whether it reacts directly to the added stimulus.
\end{enumerate}

This distinction is important because it implies that the reasoning trace is not only a bag of topics. It contains both meta-level response-format negotiation and substantive moral interpretation. The eventual coding scheme will need to preserve both layers. At the moment, no final annotation schema has been fixed, and the scheme itself remains an object of ongoing conceptual work.


\subsection{Pilot GPT-5.2 annotation of reasoning traces}

Before the fully fine-tuned MoralLens pipeline is ready, the project has already conducted a pilot annotation stage using \textbf{GPT-5.2} as an LLM judge. In this setup, the annotator receives the experiment context (system prompt, user prompt, treatment, and stimulus where available) together with the full reasoning trace and is instructed to return \emph{only} valid JSON conforming to a fixed schema. This formatting constraint is not incidental: it is essential for downstream quantitative analysis, since the pilot labels are used not only for qualitative inspection but also as inputs to statistical tests and predictive models.

The annotation prompt is written as a structured coding instruction rather than as an open-ended evaluative request. The annotator is explicitly framed as an expert in cognitive linguistics, sociology, and experimental methodology and is asked to reconstruct how the final answer was formed, which reasoning strategy was used, whether the treatment was noticed, and whether the additional stimulus entered the decision process. The instruction also contains strict parsing requirements: no markdown, no conversational preambles, and immediate output of a machine-readable JSON object.

\subsection{Pilot annotation schema}

The pilot schema currently used for reasoning traces is mixed-type and includes ordinal, binary, categorical, and multi-label fields.

\begin{itemize}
    \item \textbf{Treatment awareness} and \textbf{incentive awareness} are coded on a $0$--$3$ scale, where higher values indicate that the treatment or stimulus is more explicitly recognized and more central to the argumentation. Each score is accompanied by a short free-text explanation.
    \item \textbf{Reasoning style} includes one primary strategy (\textit{Logical-Deductive}, \textit{Emotional-Affective}, \textit{Social-Normative}, \textit{Heuristic-Shortcut}, or \textit{Identity-Based}) and a binary marker of internal conflict, indicating whether the model visibly weighs alternatives.
    \item \textbf{Epistemic certainty} is coded on a $1$--$5$ scale.
    \item \textbf{Extracted topics} are represented as a short list of key themes. These labels are included in the schema, but the present statistical analysis deliberately does not rely on topic-modeling procedures.
    \item \textbf{Social desirability} is decomposed into three separate $0$--$2$ dimensions: normative orientation, fear of judgment, and safe-answer shift.
    \item \textbf{Role adherence} includes a binary \textit{AI leakage} flag and a $1$--$3$ immersion score capturing how convincingly the model remains in the role of a human respondent.
    \item \textbf{Affective state} includes categorical labels for overall valence and dominant emotion.
    \item \textbf{Cognitive attribution} includes categorical labels for locus of control and temporal focus.
    \item \textbf{Epistemic assumptions} currently include a binary marker for hallucinated conditions, i.e.\ invented personal circumstances used to justify the answer.
\end{itemize}

This pilot schema should be viewed as an operational bridge between open-ended reasoning text and the later, more elaborate coding scheme that is still under construction. It is already rich enough to support descriptive and predictive analysis, but it remains provisional and will likely be expanded or revised once the broader coding framework is stabilized.

\subsection{Planned training-data construction}

The project plans to create an initial training dataset for reasoning annotation using a large frontier model such as DeepSeek or Gemini as an LLM judge. The choice between them has not yet been finalized. The basic plan is to use a strong model to generate candidate labels for the reasoning traces according to the emerging coding scheme, and then compare this automatic annotation against manual coding on a validation subset.

If agreement is sufficiently high, the automatically generated labels can be used as the main training signal for the annotation model. If agreement is not satisfactory, the project will move toward more complex ensemble-style annotation strategies in which multiple models generate candidate labels and disagreements are resolved through additional adjudication or manual review.

\subsection{MoralLens as a fine-tuned annotation model}

The project presentation describes MoralLens as a dedicated fine-tuned LLM for structured annotation of reasoning traces. The planned setup is instruction-following fine-tuning of a base model such as Qwen or LLaMA 3.1 8B using LoRA or QLoRA, with the full reasoning trace provided as input and a structured JSON object as output. The target is to predict membership in a psychological and cognitive taxonomy of reasoning categories, rather than to generate another free-form summary.

The presentation also specifies an anticipated training corpus of roughly 7{,}000 reasoning traces annotated automatically by a stronger model and then partially manually verified. This design is appropriate for the present project because it turns the annotation stage into a reusable measurement instrument rather than a one-off prompting heuristic.

\subsection{Planned quality criteria and evaluation metrics}

The MoralLens presentation defines several target quality criteria for the annotation pipeline.

\begin{itemize}
    \item \textbf{Inter-annotator agreement / human--LLM agreement.} The target Cohen's $\kappa$ is reported as at least $0.75$ at one point in the presentation and at least $0.80$ in the metric summary. For the coursework draft, the conservative interpretation is that agreement in this range is considered necessary before the automated labels are trusted as a training signal.
    \item \textbf{Schema validity rate.} The target schema-valid output rate is at least $0.95$, which is important because the downstream analysis requires structured JSON rather than loosely formatted explanations.
\end{itemize}

In addition, the presentation proposes task-specific predictive metrics depending on the type of coded variable:

\begin{itemize}
    \item \textbf{categorical labels:} Macro-F1;
    \item \textbf{ordinal labels:} quadratic weighted kappa (QWK) and MAE;
    \item \textbf{binary labels:} F1 and ROC-AUC;
    \item \textbf{multi-label outputs:} Micro-F1;
    \item \textbf{overall multi-category performance:} weighted Macro-F1.
\end{itemize}

These metrics are appropriate for the project because the eventual coding scheme is likely to be mixed-type rather than purely categorical. Some dimensions will be binary, some ordinal, and some multi-label.

\subsection{Why a dedicated model is preferable to ad hoc coding}

A dedicated annotation model is preferable here for three reasons. First, the reasoning traces are long and conceptually dense; they contain more than surface-level lexical cues. Second, the same coding instrument must eventually be reusable across different experiments, treatments, models, and prompt stimuli. Third, the project ultimately needs structured variables that can be used in statistical analysis, not merely impressionistic summaries. A fine-tuned model that outputs validated JSON labels is therefore a natural bridge between qualitative interpretation and quantitative inference.


\section{Preliminary Analysis of Pilot-Annotated Reasoning Traces}

\subsection{Pilot corpus and analytical status}

The current statistical analysis is based on a pilot subset of the annotated reasoning traces rather than on the final corpus. The working notebook combines four run sources: two no-stimulus baselines (\texttt{no\_p2} and \texttt{no\_t1}) and two stimulus conditions (\texttt{incentive\_C} and \texttt{incentive\_X}), yielding $5{,}600$ rows in total across four reasoning-capable models and two experiments. After dropping rows with missing markup, most predictive analyses operate on approximately $4{,}555$ complete cases.

This stage should be interpreted as a \emph{raw but already informative} pilot analysis. The same analytical toolkit is intended to be reused later on a broader and more refined reasoning annotation once that markup is ready. The immediate goal is therefore methodological: to test whether the current annotation already produces structured variables that capture stable differences between runs and are usable in statistical modeling.

\subsection{Statistical tests currently applied to the pilot markup}

The present notebook uses a layered set of statistical tests, chosen to match the scale type of the annotated variables.

For ordinal variables such as treatment awareness, incentive awareness, epistemic certainty, role immersion, and the social-desirability subscales, the baseline condition comparison is conducted using two-sided Mann--Whitney $U$ tests. Effect sizes are summarized through rank-biserial style measures derived from the test statistic. For binary variables, such as internal conflict, AI leakage, and hallucinated conditions, Pearson chi-squared tests are used together with effect sizes analogous to Cram\'er's $V$. For categorical variables such as primary reasoning strategy, affective valence, dominant emotion, locus of control, and temporal focus, chi-squared tests are likewise used to assess whether the distribution of categories differs between conditions. Because many variables are tested in parallel, the notebook applies Benjamini--Hochberg false-discovery-rate correction.

In addition to these pooled tests, the same family of nonparametric comparisons is repeated \emph{within each model} and \emph{within each experiment}. This is important because the project is explicitly comparative, and pooled significance alone could conceal heterogeneity across model families or substantive vignette domains.

Finally, for pairwise visualization of selected run-source contrasts, the notebook also computes Welch $t$-tests and Cohen's $d$ on the pilot-coded variables. These pairwise panels are not treated as the main inferential layer, but they provide a compact descriptive summary of how particular stimulus variants differ from the no-stimulus baseline and from an alternative no-stimulus run collected under a different generation setup.

\subsection{Predictive models used for the pilot markup}

Beyond hypothesis tests on individual variables, the notebook treats the annotated reasoning traces as a feature space for supervised prediction. The primary classification task is to predict whether a reasoning trace came from a stimulus condition or from a no-stimulus condition.

Several families of models are used for this purpose. First, regularized logistic regression is fitted as a baseline linear classifier, with cross-validated ROC-AUC as the main evaluation metric. Second, tree-based models are used to capture nonlinearity and interactions, specifically random forests and gradient boosting classifiers. Their performance is likewise evaluated by cross-validated ROC-AUC. In the tree-based setting, permutation importance is used to identify which markup features are most consequential for discriminating stimulus from baseline runs.

The notebook also fits these classifiers at a more granular level. There are models estimated separately by \texttt{(model, experiment)} cells, and there are pairwise run-source classifiers contrasting \texttt{no\_p2} with \texttt{incentive\_C}, \texttt{no\_p2} with \texttt{incentive\_X}, and \texttt{no\_p2} with \texttt{no\_t1}. This last comparison is particularly important because it provides a rough benchmark for comparing \emph{stimulus-induced separability} against separability induced by differences in generation setup unrelated to the stimulus itself.

In parallel, the notebook estimates ordered-logit models for ordinal markup outcomes and simplified mixed-effects models in which the condition predicts a given annotation variable while controlling for model and experiment. These models are not yet the final inferential framework of the project, but they already serve an important role: they show which parts of the markup move systematically with stimulus assignment after at least partial adjustment for model family and vignette domain.

\subsection{Preliminary results from the pilot markup}

The main pilot finding is that reasoning traces are clearly distinguishable between the stimulus and no-stimulus setups, and this separability is substantially stronger than the separability observed between two no-stimulus runs collected under different generation settings. In the pooled condition-classification task, cross-validated ROC-AUC is approximately $0.89$ for logistic regression, $0.91$ for random forest, and $0.94$ for gradient boosting. At the finer \texttt{model $\times$ experiment} level, the classifiers remain strong in most strata, typically in the $0.88$--$0.95$ ROC-AUC range.

The most informative comparison is the pairwise one. For the \texttt{no\_p2} versus \texttt{no\_t1} contrast---that is, a no-stimulus versus no-stimulus comparison reflecting a different generation setup rather than a different realism stimulus---the cross-validated AUC is modest for the models with balanced data, roughly $0.61$ for Kimi and $0.67$ for GLM under logistic regression. By contrast, the \texttt{no\_p2} versus \texttt{incentive\_C} contrast yields much stronger discrimination, around $0.96$ for Kimi and $0.97$ for GLM, with similarly high performance for random forests. The \texttt{no\_p2} versus \texttt{incentive\_X} contrast is also clearly separable, with AUC values around $0.88$--$0.94$ depending on model and classifier. In substantive terms, this means that the reasoning markup captures a stronger shift between stimulus and baseline conditions than between baseline runs that differ only in other technical settings. This is precisely the pattern one would expect if the stimulus is entering the reasoning process in a meaningful way.

Several variable-level differences also emerge in the pilot stage. In pooled comparisons, stimulus conditions are associated with substantially lower \textit{social desirability normative orientation} and lower \textit{social desirability total}, while \textit{incentive awareness} increases sharply by construction and \textit{AI leakage} becomes more frequent. Distributional tests also indicate that primary reasoning strategy, affective valence, dominant emotion, locus of control, and temporal focus differ significantly between stimulus and no-stimulus conditions after multiple-testing correction. These findings are still preliminary and partly shaped by the current markup design, but they suggest that the stimuli affect not only the final answer but also the argumentative and stylistic organization of the reasoning trace.

\subsection{Why this pilot matters for the next annotation stage}

Methodologically, the pilot analysis serves as a proof of concept for the broader reasoning-analysis program. It shows that the current annotation prompt already yields variables with enough structure to support nontrivial statistical testing and supervised prediction. Just as importantly, it provides a template for the next annotation stage: the same combination of nonparametric tests, regression-style models, mixed-effects comparisons, and condition classifiers can later be applied to a richer coding scheme once the next markup is available.

The present pilot results should therefore not be read as definitive substantive conclusions. Their value lies in demonstrating that the markup is analytically live: it already captures stimulus-related structure, and it appears to do so more strongly than it captures variation associated with a mere change in temperature or related generation settings.

\section{Interpretable AI Extensions}

A further line of work is planned at the interface between reasoning annotation and interpretable AI. The idea is to use the structured markup features not only as descriptive variables but also as explanatory inputs for models that predict the final survey answer. In this setup, the final answer becomes the target, and the annotated reasoning features become candidate explanatory factors.

At the current stage, this block remains intentionally modest and largely prospective. The main tools under consideration are model-agnostic explanation methods such as SHAP and LIME, as well as simpler coefficient-based interpretations for sparse linear models. The goal is to estimate which reasoning features most strongly push the model toward one final answer rather than another and whether those local or global importance patterns differ across prompt stimuli, experiments, and model families.

This future block is especially important because it provides a bridge between descriptive markup and mechanism-oriented explanation. Rather than merely saying that stimulus and no-stimulus traces differ, the project will be able to ask \emph{which annotated properties of the reasoning} account for the final answer and whether those properties remain stable across alternative coding schemes. For that reason, the same interpretable-model toolkit is intended to be reused not only on the current pilot markup but also on the next, more developed markup that is still under preparation.

\section{Hidden-State Analysis}

In addition to the text-level analysis of reasoning traces, the project also includes a representational analysis of hidden states. The goal of this component is to test whether prompt stimuli alter not only the textual output of the model but also its internal representations, and if so, at what depth in the network this becomes visible.

\subsection{Layer-wise stimulus separability}

The first question is whether hidden states can distinguish between setups \emph{with} and \emph{without} a stimulus.

For each layer $k$, the project extracts a hidden-state representation and trains a regularized linear probe of the form
\[
\texttt{stimulus\_present} \sim \texttt{hidden\_state\_layer\_k}.
\]

The main object of interest is not the coefficient vector itself but the predictive performance of the probe on held-out data. The primary evaluation criteria are expected to include accuracy, ROC-AUC where appropriate, and comparison against a permutation baseline.

The substantive question is whether the effect of the stimulus emerges in the model's internal representations and, if it does, at which layers it becomes linearly decodable. This would provide evidence that the system-prompt intervention is not merely a superficial textual artifact but enters the model's internal processing trajectory.

\subsection{Layer-wise answer predictability}

The second question is whether hidden states at different layers can predict the model's eventual final answer.

For each layer $k$, the project trains a probe of the form
\[
\texttt{final\_answer} \sim \texttt{hidden\_state\_layer\_k}.
\]

This analysis is run separately for the no-stimulus and stimulus conditions. The central comparison is therefore not just where the answer becomes decodable, but whether the presence of a stimulus changes the timing or strength of that decodability.

The main substantive question is whether the stimulus makes the eventual answer visible \emph{earlier}, \emph{later}, or \emph{more strongly} in the internal states compared with baseline. If answer predictability shifts systematically across layers under stimulus conditions, this would suggest that the prompt intervention is altering the trajectory of decision formation rather than only the wording of the final response.

\subsection{Why this matters for the project}

The hidden-state analysis complements the textual reasoning analysis in an important way. Reasoning traces are observable and interpretable, but they are also generated narratives, and their relation to the model's internal computation is not trivial. Probe-based hidden-state analysis does not solve that interpretability problem completely, but it offers a second empirical window. If both the textual reasoning trace and the hidden-state probes indicate a systematic difference between stimulus and no-stimulus conditions, the case for a real intervention effect becomes stronger.

\section{Planned Analytical Strategy}

The present coursework draft does not yet report final empirical results, but the project already has a sufficiently articulated empirical design to be described as a coherent research pipeline. The purpose of this section is therefore not only to list the statistical procedures planned for later stages, but to summarize the full logic of the study in the order in which its components are expected to work together. In this sense, the analytical strategy is inseparable from the broader research design: the project begins with the construction of an experimental corpus, proceeds through controlled LLM data generation, then moves to several layers of analysis of the generated material, and finally integrates these layers into a single interpretation of how LLMs respond to vignette-based stimuli in socially sensitive contexts.

At the most general level, the project studies LLM behavior in vignette-based survey experiments on social desirability bias by observing it on four connected levels: final answers, verbalized reasoning traces, structured annotations of those reasoning traces, and internal hidden-state representations. The central question is not only whether a model reproduces a known survey effect, but also how that effect is organized, where in the generation process it becomes visible, and how it changes under different prompt-based realism stimuli.

\subsection{Stage 1: corpus construction and experimental scope}

The first stage of the project is the construction of a corpus of survey and vignette experiments suitable for LLM-based replication. Conceptually, the broader research agenda distinguishes three types of SDB-related experimental settings: method-oriented designs, self-report designs, and attitude-expression designs. The current coursework focuses on the third category, since attitude-expression experiments are both substantively central to the problem of social desirability and relatively tractable for a first stage of synthetic piloting.

At the present stage, the empirical core of the project consists of two experiments already adapted for model-based replication. The first is a negative case centered on support for torture or coercive interrogation; the second is a positive case centered on willingness to donate. This pairing is analytically important because it allows the project to compare model behavior in settings where the socially undesirable response is associated with approval of normatively condemned conduct and settings where the socially desirable response is associated with prosocial self-presentation. In other words, the project is designed from the beginning to cover both socially negative and socially positive domains rather than only one type of morally charged response.

\subsection{Stage 2: controlled data generation with multiple reasoning models}

The second stage is the controlled collection of synthetic responses from several reasoning-capable LLMs. The current model set includes Kimi K2 Thinking, Qwen3-9B, GLM-4.6, and MiniMax-M2.5. These models were selected not as a representative sample of all available LLMs, but as a practical set of contemporary reasoning-oriented systems differing in family, scale, and alignment profile. This makes it possible to treat cross-model comparison as a substantive part of the design rather than as a robustness check added at the end.

For each run, the model is exposed to the experimental material and produces, in a single generation pass, two outputs that are both preserved for analysis: the final answer required by the survey instruction and the reasoning trace leading to that answer. The project does not rely on a two-stage setup in which reasoning and answer are elicited separately; instead, it records them jointly under the same prompt conditions. Temperature is fixed at \(0.00\) in the main design in order to maximize reproducibility and reduce uncontrolled variability in the collected sample. Thus, the basic observation unit of the project is a model response consisting of a final outcome and a reasoning trace generated under a known combination of experiment, treatment condition, model family, and system-prompt stimulus.

\subsection{Stage 3: outcome-level replication of vignette effects}

The first analytical layer concerns the observable survey outcome. Here the goal is to determine whether the models reproduce the qualitative direction of the original vignette effects within each study. In the torture case, this means asking whether changes in the scenario alter support for coercive interrogation in the expected direction. In the donation case, it means asking whether the framing of the situation shifts willingness to donate in the expected way. At this stage, the primary benchmark is directional replication rather than precise recovery of published coefficients or exact effect sizes.

Empirically, this layer compares response distributions across experimental conditions, models, and stimuli. Depending on the scale of the dependent variable, this can involve descriptive contrasts, contingency-table analysis, chi-squared tests, proportion tests, and regression-style specifications in which the final answer is modeled as a function of treatment, stimulus, model family, and their interactions. In simplified form, the main specification can be written as
\[
\Pr(Y_i = 1) = \operatorname{logit}^{-1}\!\left(\beta_0 + \beta_1 \texttt{Treatment}_i + \beta_2 \texttt{Stimulus}_i + \beta_3 \texttt{Model}_i + \beta_4 (\texttt{Treatment}_i \times \texttt{Stimulus}_i)\right),
\]
where \(Y_i\) denotes the final answer for observation \(i\). Since the project is conceived as synthetic piloting rather than exact statistical replication of human samples, the main substantive question is whether the structure of the original contrast survives in model-generated data.

\subsection{Stage 4: system-prompt stimuli as an experimental factor}

The second analytical layer concerns the additional realism-oriented system prompts. These prompts are not treated as a miscellaneous collection of ad hoc manipulations. Rather, they form a structured design space. Across the current stimulus inventory, the project varies at least three substantively meaningful dimensions: the harshness of the realism instruction, the presence or absence of an anti-overcorrection clause, the identity of the actor who bears the consequences of unrealistic responses, and the valence of the message, that is, whether realism is framed through loss or through gain.

This design allows the project to ask not only whether a given named prompt works, but what properties of prompts are actually responsible for behavioral change. At the present pilot stage, three provisional expectations already guide the analysis. First, prompts that are too harsh and insufficiently balanced may induce sharper or more extreme answers. Second, positively framed realism prompts may be more effective than negatively framed ones in eliciting socially undesirable responses. Third, changing the nominal recipient of the consequences may matter less than expected, at least within the negatively framed variants. These observations are still preliminary, but they are already important because they define stimulus analysis as a genuine explanatory part of the project rather than a superficial prompt comparison exercise.

\subsection{Stage 5: reasoning traces as data in their own right}

The third analytical layer concerns the reasoning traces themselves. The project does not treat them as decorative output or as a mere transparency artifact. Instead, it assumes that reasoning traces are a separate empirical object that may reveal distinctions invisible at the level of the final answer alone. Two identical answers may be reached through different argumentative paths; conversely, the same treatment effect may appear stable in aggregate outcomes while the underlying reasoning shifts substantially across stimuli or across models.

At the current stage, exploratory qualitative reading suggests that relevant reasoning features fall into at least two broad families. The first family contains format-related features, such as explicit reflection on scale use, reluctance to choose an extreme category, concern with sounding too virtuous or too harsh, and differentiation between what a respondent would think and what a respondent would actually do. The second family contains content-related features, such as recognition of where the treatment is embedded in the vignette, appeal to moral norms, references to consequences, sensitivity to social acceptability, and direct reaction to the realism stimulus. A finalized full coding scheme is still under development, but the project already treats the conversion of reasoning traces into structured variables as one of its central methodological tasks.

\subsection{Stage 6: pilot markup of reasoning traces}

Before building the final annotation system, the project has already conducted a pilot markup of part of the reasoning corpus using GPT-5.2 as an external coding instrument. In this setup, GPT-5.2 does not act as a respondent. Instead, it receives the experiment context together with the full reasoning trace and returns a structured JSON annotation according to a predefined instruction. This pilot schema includes such dimensions as treatment awareness, incentive awareness, primary reasoning style, internal conflict, epistemic certainty, extracted topics, several components of social desirability, role adherence, affective state, cognitive attribution, and epistemic assumptions.

This pilot stage is important for two reasons. First, it tests whether reasoning traces can already be transformed into variables suitable for systematic analysis. Second, it provides an empirical basis for refining the future coding scheme. In other words, the project is not waiting for the final annotation taxonomy before beginning quantitative analysis; instead, it uses pilot markup as an intermediate methodological layer that helps clarify what kinds of regularities are already detectable in the reasoning corpus.

\subsection{Stage 7: statistical analysis of the pilot markup}

A separate block of analysis is devoted specifically to the already marked-up reasoning traces. This analysis deliberately excludes topic modeling or STM and instead focuses on statistical tests and predictive models. Depending on variable type, the current workflow includes cross-tabulation, chi-squared tests for categorical features, Mann--Whitney comparisons for ordinal scores, effect-size summaries, and supervised classification models that attempt to distinguish conditions from markup-derived features alone.

The logic of this stage is straightforward: if the markup is informative, then reasoning-derived features should help separate traces produced under different conditions. In simplified form, the current classification setup can be written as
\[
\Pr(\texttt{StimulusPresent}_i = 1) = \operatorname{logit}^{-1}\!\left(\alpha_0 + \sum_{j=1}^{p}\alpha_j X_{ij}\right),
\]
where \(X_{ij}\) are annotated reasoning features for observation \(i\). The main preliminary finding at this stage is that reasoning traces from runs with and without the realism stimulus appear more separable than traces generated under different temperature settings. Methodologically, this is a promising result: it suggests that the stimulus affects the structure of reasoning more strongly than moderate variation in decoding parameters does. The same family of tests and predictive models is planned for reuse later, once the richer coding scheme is completed.

\subsection{Stage 8: from markup features to answer-level mechanisms}

Once reasoning traces are translated into structured variables, the next analytical step is to relate those variables to the final answer. This stage is designed to move from descriptive pattern detection to mechanism-oriented analysis. Substantively, it asks whether specific reasoning features predict socially desirable versus socially undesirable outcomes, whether particular stimuli increase treatment awareness or moral hedging, and whether observed changes in answers are statistically associated with changes in reasoning structure.

This stage will rely on regression-style models in which the dependent variable may be either the final answer or a reasoning feature, while the predictors may include treatment condition, stimulus type, model family, and interactions among them. A simplified example is
\[
\Pr(Y_i = 1) = \operatorname{logit}^{-1}\!\left(\gamma_0 + \sum_{j=1}^{p}\gamma_j X_{ij} + \gamma_{p+1}\texttt{Treatment}_i + \gamma_{p+2}\texttt{Stimulus}_i + \gamma_{p+3}\texttt{Model}_i \right),
\]
where \(X_{ij}\) are markup-derived reasoning features. This block is central because it connects variation in survey answers to interpretable variation in argumentative structure.

\subsection{Stage 9: future annotation model and scalable reasoning analysis}

The pilot markup is not the final intended endpoint. The larger plan is to create a more stable and scalable annotation pipeline by constructing a training dataset, comparing machine markup with human verification, and then fine-tuning a dedicated model for reasoning annotation in the spirit of the MoralLens tool. At that stage, the project will move from pilot GPT-based markup to a reusable annotation instrument capable of processing larger reasoning corpora more consistently and more cheaply.

This part of the design is still under construction, especially with regard to the final coding taxonomy and labeling workflow. However, even in its current form the plan is already clear: reasoning traces will eventually be analyzed not only through ad hoc pilot markup, but through a dedicated annotation model whose quality will itself be evaluated through agreement and predictive metrics.

\subsection{Stage 10: interpretable-AI layer}

A further planned extension is to apply interpretable-AI methods to models linking annotated reasoning features with final answers. Once a sufficiently stable feature space is available, predictors such as logistic regression or other supervised classifiers can be supplemented with tools such as SHAP or LIME in order to identify which reasoning features contribute most strongly to a predicted answer. In simplified notation, for a fitted predictor \(f(x)\), the goal is to decompose a prediction into feature-level contributions,
\[
f(x) \approx \phi_0 + \sum_{j=1}^{p}\phi_j,
\]
where \(\phi_j\) denotes the contribution of feature \(j\). At the present stage, this layer remains mostly prospective, but it is already part of the overall design because it provides a bridge between statistical association and interpretable explanation.

\subsection{Stage 11: internal-representation analysis}

The final analytical block moves below the text level and examines internal model representations. In addition to collecting final answers and reasoning traces, the project plans to run direct inference and compare hidden states between baseline prompts and prompts augmented with realism stimuli. This is meant to test whether stimulus effects are visible in the internal dynamics of the model even before they become clearly visible in the final answer.

Two probe-based analyses are planned. The first is \emph{layer-wise stimulus separability}. For each layer, a regularized linear probe will be trained to predict whether the prompt included a stimulus:
\[
\texttt{stimulus\_present} \sim \texttt{hidden\_state}_{k}.
\]
The key output here is not the probe coefficients themselves but held-out predictive quality, such as accuracy, ROC-AUC, and comparison against a permutation baseline. The second is \emph{layer-wise answer predictability}. For each layer, a probe will be trained to predict the model's eventual final answer:
\[
\texttt{final\_answer} \sim \texttt{hidden\_state}_{k}.
\]
This second analysis will be conducted separately for baseline and stimulus conditions. The substantive question is whether the realism stimulus shifts the point at which the future answer becomes decodable from internal states.

\subsection{Integrative interpretation}

Taken together, these stages form a single empirical plan rather than a loose collection of possible analyses. The project begins by selecting suitable SDB-related vignette experiments, then generates synthetic responses from several reasoning-capable LLMs under controlled prompt conditions, then analyzes observable answer patterns, then studies reasoning traces both qualitatively and through pilot structured markup, then links markup-derived features back to final answers, and finally extends the analysis both downward to hidden states and outward to interpretable-AI tools. Some of these stages have already been implemented in pilot form, while others remain planned for later development. Nevertheless, the full architecture of the research is already visible. The broader aim is therefore not only to test whether LLMs can imitate survey respondents, but to understand under what conditions they can serve as a credible instrument for piloting vignette-based experimental designs in domains where social desirability bias is especially salient.

\section{Scope and Current Boundaries}

Several limits of the present draft should be made explicit.

\begin{itemize}
    \item The current empirical corpus is intentionally narrow and includes only two attitude-expression experiments.
    \item The coding scheme for reasoning traces is still under construction; the project currently has a distinction between format-related and substantive features, but not yet a frozen final taxonomy.
    \item The automated annotation pipeline is planned in detail, but the final training model and the exact supervising model for initial auto-labeling have not yet been fixed.
    \item The hidden-state analysis has a clear design logic but remains a planned component rather than a completed empirical section.
    \item The draft does not present final estimates or conclusions; it defines the architecture of the study.
\end{itemize}

These limitations are acceptable for the current stage because the purpose of the text is to provide a coherent and extensible research design for subsequent empirical implementation.

\section{Conclusion}

This project studies synthetic sample generation in the context of experimental research by focusing on a demanding subclass of vignette-based survey experiments: cases shaped by social desirability bias. Its central question is how large language models react to stimuli in such experiments. The design begins with two attitude-expression cases that capture opposite normative directions, one negative and one positive, and examines how system-prompt interventions modify the final answer, the reasoning trace, and eventually the model's internal representations.

Methodologically, the project combines three elements: reproducible prompt-based experimentation, structured annotation of reasoning traces through the developing MoralLens instrument, and layer-wise analysis of hidden states. Together, these components are intended to move beyond the narrow question of whether a model can output a plausible answer. The deeper goal is to understand whether LLMs can function as informative pilot environments for experimental survey design precisely in domains where socially desirable responding is most likely to distort both human and model behavior.

\end{document}
